% 四种多面体：(a)空集 x>=0, x1+x2<=1, x1+x2>=2；(b)无界区域 x>=0, x1+x2>=1 上 min x1+2x2 在(1,0)达到；
% (c)同一区域上 min x2-x1 沿射线(1,0)+t(1,0)无下界；(d)条带 0<=x2<=1（x1自由）包含直线、没有极点。
\begin{tikzpicture}[x=8.5mm,y=8.5mm,
  note/.style={font=\FigureNoteFont\fontsize{7.5}{9}\selectfont,text=FigureMuted},
  ttl/.style={font=\fontsize{8}{10}\selectfont,anchor=north}]
  % (a)
  \begin{scope}
    \fill[FigureBlueFill] (0,0)--(1,0)--(0,1)--cycle;
    \fill[FigureGreenFill] (2,0)--(3,0)--(3,2.6)--(0,2.6)--(0,2)--cycle;
    \draw[FigureBlue,line width=.9pt] (1,0)--(0,1);
    \draw[FigureGreen,line width=.9pt] (2,0)--(0,2);
    \draw[->,FigureStroke,line width=.6pt] (-0.2,0)--(3.2,0) node[right]{$x_1$};
    \draw[->,FigureStroke,line width=.6pt] (0,-0.2)--(0,2.8) node[above]{$x_2$};
    \node[note,anchor=north] at (0.6,-0.02) {$x_1+x_2\le1$};
    \node[note] at (1.9,1.7) {$x_1+x_2\ge2$};
    \node[ttl] at (1.5,-0.45) {(a) 空集};
  \end{scope}
  % (b)
  \begin{scope}[xshift=40mm]
    \fill[FigureBlueFill] (1,0)--(3,0)--(3,2.6)--(0,2.6)--(0,1)--cycle;
    \draw[FigureBlue,line width=.9pt] (3,0)--(1,0)--(0,1)--(0,2.6);
    \begin{scope}\clip (0,0) rectangle (3,2.6);
      \foreach \k in {2,3}{\draw[FigureGreen,line width=.6pt,dashed] (0,\k/2)--(\k,0);}
      \draw[FigureGreen,line width=1.1pt] (0,0.5)--(1,0);
    \end{scope}
    \draw[->,FigureStroke,line width=.6pt] (-0.2,0)--(3.2,0) node[right]{$x_1$};
    \draw[->,FigureStroke,line width=.6pt] (0,-0.2)--(0,2.8) node[above]{$x_2$};
    \fill[FigureRed] (1,0) circle (1.6pt);
    \draw[->,FigureRed,line width=.9pt] (2.2,1.6)--++(-0.27,-0.54);
    \node[note,anchor=south] at (2.2,1.62) {$-(1,2)$};
    \node[note,anchor=north] at (1,-0.02) {$(1,0)$};
    \node[ttl] at (1.5,-0.45) {(b) 无界，目标有下界};
  \end{scope}
  % (c)
  \begin{scope}[xshift=80mm]
    \fill[FigureBlueFill] (1,0)--(3,0)--(3,2.6)--(0,2.6)--(0,1)--cycle;
    \draw[FigureBlue,line width=.9pt] (3,0)--(1,0)--(0,1)--(0,2.6);
    \begin{scope}\clip (0,0) rectangle (3,2.6);
      \foreach \k in {-2,-1,0,1}{\draw[FigureGreen,line width=.6pt,dashed] (0,\k)--(4,\k+4);}
    \end{scope}
    \draw[->,FigureStroke,line width=.6pt] (-0.2,0)--(3.2,0) node[right]{$x_1$};
    \draw[->,FigureStroke,line width=.6pt] (0,-0.2)--(0,2.8) node[above]{$x_2$};
    \draw[->,FigureRed,line width=1.1pt] (1,0.08)--(2.85,0.08);
    \fill[FigureRed] (1,0) circle (1.6pt);
    \node[note,anchor=north] at (1,-0.02) {$\bm x_0$};
    \node[note,anchor=south] at (2.05,0.12) {$\bm x_0+t\bm d$};
    \node[ttl] at (1.5,-0.45) {(c) 无界，目标无下界};
  \end{scope}
  % (d)
  \begin{scope}[xshift=120mm]
    \fill[FigureBlueFill] (-0.4,0)--(3,0)--(3,1)--(-0.4,1)--cycle;
    \draw[FigureBlue,line width=.9pt] (-0.4,0)--(3,0) (-0.4,1)--(3,1);
    \draw[<->,FigureRed,line width=.9pt] (-0.3,0.5)--(2.9,0.5);
    \draw[->,FigureStroke,line width=.6pt] (-0.4,0)--(3.2,0) node[right]{$x_1$};
    \draw[->,FigureStroke,line width=.6pt] (0,-0.2)--(0,2.8) node[above]{$x_2$};
    \node[note] at (1.5,1.55) {$0\le x_2\le1$};
    \node[note] at (1.5,1.95) {含整条直线};
    \node[ttl] at (1.5,-0.45) {(d) 条带，没有极点};
  \end{scope}
\end{tikzpicture}
